% Chapter 2 --- Measurable cones (paper §3, milestone M2).
%
% Mirrors theories/mcones/:
%   ar.v          paper §3 — the $\Ar$ parameter (a small full subcategory
%                            of $\Meas$ containing 0 and closed under
%                            products), plus product helpers.
%   mcone.v       paper Defs 3.2 / 3.6 / 3.7 — tests, measurable cones,
%                            measurable paths, Lemmas 3.9 / 3.10.
%   fmeas.v       paper §3.2.1 — finite measures $\FMeas(X)$, with cone
%                            structure, $\MCone$ structure, and the
%                            pushforward functor (Lemma 3.17).
%   path.v        paper §3.2.2 — paths $\Path(X, B)$ as an $\MCone$, plus
%                            the flattening iso (Lemma 3.19).
%   mcone_cat.v   paper Prop 3.11 / Defs 3.13 / 3.15 — the category
%                            $\MCones$ and the rescaling $\alpha_B$.

% Local macros used only by this chapter.
\providecommand{\mzero}{\mathbf{0}}
\providecommand{\arzero}{\mathbf{0}}
\providecommand{\testof}[2]{\mathrm{test}_{#1}(#2)}
\providecommand{\reindex}[2]{#2 \circ #1}
\providecommand{\evtest}[1]{e_{#1}}
\providecommand{\trianglerhd}{\mathbin{\triangleright}}

%%%%%%%%%%%%%%%%%%%%%%%%%%%%%%%%%%%%%%%%%%%%%%%%%%%%%%%%%%%%%%%%%%%%%%%%
\chapter{Measurable cones}\label{ch:mcones}

This chapter formalises paper \S3: the \emph{measurability structure}
on a cone, the category $\MCones$, and the two basic examples of
measurable cones --- the cone of finite measures $\FMeas(X)$ and the
cone of paths $\Path(X, B)$. All material lives under
\texttt{theories/mcones/}.

The leading idea is that a cone $P$ becomes \emph{measurable} when it
is equipped with a small family $\Mtest = (\Mtest_X)_{X \in \Ar}$ of
$[0,1]$-valued \emph{tests}, indexed by a parameter category $\Ar$ of
measurable spaces. The tests play the role of a separating set of
bounded linear functionals; in the discrete-cones world they
correspond to dual-cone elements (\S\ref{sec:mcones-def} below). This
is the cone-theoretic analogue of a quasi-Borel space structure: paths
into the cone are detected by composing with tests, and a function
$\gamma : X \to P$ is a \emph{measurable path} exactly when every test
``sees'' it as a measurable real-valued function.

%%%%%%%%%%%%%%%%%%%%%%%%%%%%%%%%%%%%%%%%%%%%%%%%%%%%%%%%%%%%%%%%%%%%%%%%
\section{The parameter \texorpdfstring{$\Ar$}{Ar}}\label{sec:ar}

\paragraph{Paper setup.}
The whole of \S3 is parametrised by a fixed small full subcategory
$\Ar$ of the category $\Meas$ of measurable spaces, with the following
two closure properties:
\begin{itemize}
  \item $\Ar$ contains the terminal object $\arzero$ of $\Meas$
    (the one-point measurable space, written $\arzero$ in the paper
    because it is ``geometrically $0$-dimensional'');
  \item $\Ar$ is closed under finite cartesian products.
\end{itemize}
All objects of $\Ar$ are required to be non-empty. Morphisms in $\Ar$
are measurable functions. Remark~3.1 of the paper notes that
$\Ar = \{ \mathbb{R}^n \mid n \ge 0 \}$ is the canonical choice; we
keep $\Ar$ a parameter for flexibility (see the design notes below).

\begin{definition}[Paper \S3, the $\Ar$ parameter]\label{def:meas-subcat}
  \uses{}
  \rocq{Icones.mcones.ar.MeasSubcat}\rocqok
  An $\Ar$-datum is a record \texttt{MeasSubcat R} consisting of:
  \begin{itemize}
    \item a type \texttt{ar\_obj} of object names;
    \item for each $X \in \Ar$, a measurable-space carrier
      $\ar\_carrier X$ with a canonical inhabitant
      $\ar\_point X : \ar\_carrier X$ (the paper's non-emptiness
      assumption, made constructive);
    \item a distinguished object $\arzero \in \Ar$ whose carrier is a
      singleton ($\ar\_zero\_singleton$ asserts $\forall x, y, x = y$);
    \item a binary product
      $\ar\_prod : \ar\_obj \to \ar\_obj \to \ar\_obj$ together with
      \emph{propositional} equalities
      $\ar\_prod\_disp\_eq$ and $\ar\_prod\_carrier\_eq$ identifying the
      carrier of $\ar\_prod\ X\ Y$ with the product
      $(\ar\_carrier\ X \times \ar\_carrier\ Y)$ at the type level, and
      two cast-measurability fields
      $\ar\_prod\_cast\_meas$ /
      $\ar\_prod\_uncast\_meas$ asserting that the two transports
      across this equality are measurable.
  \end{itemize}
  Morphisms in $\Ar$ are measurable functions, encoded via
  \mathcompA's bundle \texttt{\{mfun aT >-> rT\}}, abbreviated
  \texttt{ar\_hom Ar Y X}.
\end{definition}

The cast-measurability fields require comment: the propositional
equality $\ar\_prod\_carrier\_eq$ identifies the two carriers at type
\texttt{Type} but does not by itself relate their $\sigma$-algebras,
hence the explicit measurability obligations on the two
\texttt{eq\_rect} casts. In every standard realisation
(where $\ar\_prod\ X\ Y$ \emph{is} the mathcomp-analysis product
measurable space, so the equality is $\mathsf{eq\_refl}$), both fields
are discharged by $\measurable\_id$.

\begin{lemma}[Universal property of $\ar\_prod$]\label{lem:ar-prod-up}
  \uses{def:meas-subcat}
  \rocq{Icones.mcones.ar.ar_pair_fst}\rocqok
  \rocq{Icones.mcones.ar.ar_pair_snd}\rocqok
  \rocq{Icones.mcones.ar.ar_pair_uniqueE}\rocqok
  The maps $\ar\_prod\_fst : \ar\_prod\ X\ Y \to X$ and
  $\ar\_prod\_snd : \ar\_prod\ X\ Y \to Y$ are measurable, and the
  pairing $\ar\_pair\ f\ g$ of $f \in \Ar(Z, X)$ and $g \in \Ar(Z, Y)$
  is the unique $\Ar$-morphism through which $f$ and $g$ factor.
\end{lemma}
\begin{proof}
  The projection $\ar\_prod\_fst$ is the composite of
  $\mathsf{fst}$ with the type-level uncast, whose measurability is
  the field $\ar\_prod\_uncast\_meas$. Symmetrically for the second
  projection. Pairing $\ar\_pair\ f\ g$ is the composite of the
  bundled pair $z \mapsto (f\ z,\ g\ z)$ with the cast in the reverse
  direction. Uniqueness reduces to the cast/uncast round-trip
  identities $\ar\_prod\_castK$ and $\ar\_prod\_uncastK$.
\end{proof}

%%%%%%%%%%%%%%%%%%%%%%%%%%%%%%%%%%%%%%%%%%%%%%%%%%%%%%%%%%%%%%%%%%%%%%%%
\section{Measurability structures and measurable
  cones}\label{sec:mcones-def}

%% Tests at an arity X. ------------------------------------------------

\begin{definition}[Paper Def 3.2, the tests of arity $X$]%
  \label{def:test-of}
  \uses{def:meas-subcat}
  \rocq{Icones.mcones.mcone.test_of}\rocqok
  Let $\Ar$ be as in Definition~\ref{def:meas-subcat}, $X \in \Ar$ and
  $C$ a cone. A \emph{test of arity $X$ on $C$} is a function
  \[
    m : \ar\_carrier X \to C \to \R
  \]
  satisfying, in the Rocq record \texttt{test\_of\ Ar\ X\ C}:
  \begin{itemize}
    \item ($m \ge 0$): $m\ r\ x \ge 0$ for every $r, x$
      (\rocq{Icones.mcones.mcone.test_ge0});
    \item ($\mathrm{Msmeas}$, measurability): for every $x \in B_C$
      (i.e.\ $\cnorm{x} \le 1$), the partial application
      $r \mapsto m\ r\ x$ is a measurable function
      $\ar\_carrier X \to \R$
      (\rocq{Icones.mcones.mcone.test_meas});
    \item ($\mathrm{Msmeas}$, boundedness): $m\ r\ x \le 1$ whenever
      $\cnorm{x} \le 1$ (\rocq{Icones.mcones.mcone.test_le1});
    \item linearity in the second argument: $m\ r\ 0 = 0$,
      $m\ r\ (x + y) = m\ r\ x + m\ r\ y$, and
      $m\ r\ (s \cdot x) = s \cdot m\ r\ x$ for $s \in \Rnn$
      (\rocq{Icones.mcones.mcone.test_lin0},
      \rocq{Icones.mcones.mcone.test_linD},
      \rocq{Icones.mcones.mcone.test_linZ});
    \item $\omega$-continuity in the second argument: for every
      increasing chain $(u_n)$ bounded in $B_C$ and every upper bound
      $N$ of $(m\ r\ u_n)$, $m\ r\ (\bigsup_n u_n) \le N$
      (\rocq{Icones.mcones.mcone.test_cont});
    \item operator-norm bound: $m\ r\ x \le \cnorm{x}$ for every
      $r, x$ (\rocq{Icones.mcones.mcone.test_norm_le}).
  \end{itemize}
\end{definition}

\begin{remark}
  The Rocq encoding of a test bundles the measurability,
  non-negativity, linearity, $\omega$-continuity and norm-bound
  proofs in a single \texttt{Record}; this is the cone-side
  counterpart of mathcomp-analysis's \texttt{\{mfun \_ >-> \_\}}
  bundle, but with codomain $\R$ rather than a general
  $\mathsf{measurableType}$. Equality of two tests reduces to
  pointwise equality of the underlying functions
  (\rocq{Icones.mcones.mcone.test_eq}, by $\Prop$-irrelevance).
\end{remark}

%% Reindexing of a test. ----------------------------------------------

\begin{definition}[Paper Def 3.2, reindexing $m \circ (\varphi \times C)$]
  \label{def:test-reindex}
  \uses{def:test-of}
  \rocq{Icones.mcones.mcone.test_reindex}\rocqok
  Given $\varphi \in \Ar(Y, X)$ and a test $m$ of arity $X$, the
  \emph{reindexed test} $\reindex{\varphi}{m}$ of arity $Y$ is the
  function $(s, x) \mapsto m\ (\varphi\ s)\ x$. The seven
  test axioms transport along $\varphi$ because $\varphi$ is itself
  measurable.
\end{definition}

%% The four axioms of a measurability structure. ----------------------

\begin{definition}[Paper Def 3.2, a measurability structure]%
  \label{def:isMCone}
  \uses{def:test-of, def:test-reindex}
  \rocq{Icones.mcones.mcone.isMCone}\rocqok
  Let $C$ be a cone. A \emph{measurability structure} on $C$ is the
  data of a family
  \[
    \Mtest = \bigl(\Mtest_X\bigr)_{X \in \Ar},
    \qquad
    \Mtest_X \subseteq \testof{X}{C}
  \]
  of subsets of the tests of arity $X$, subject to the four axioms:
  \begin{description}
    \item[$(\mathrm{Msmeas})$] built into Definition~\ref{def:test-of}:
      each $m \in \Mtest_X$ is jointly measurable in its first
      argument on the unit ball, $[0,1]$-valued there, $\Rnn$-linear
      and $\omega$-continuous in its second argument.
    \item[$(\mathrm{Mscomp})$, closure under reindexing] for every
      $\varphi \in \Ar(Y, X)$ and $m \in \Mtest_X$, the reindexed test
      $\reindex{\varphi}{m}$ is in $\Mtest_Y$
      (\rocq{Icones.mcones.mcone.isMCone.mcone_M_comp}).
    \item[$(\mathrm{Mssep})$, separation] if $x_1, x_2 \in C$ satisfy
      $m\ \arzero\ x_1 = m\ \arzero\ x_2$ for every
      $m \in \Mtest_{\arzero}$, then $x_1 = x_2$
      (\rocq{Icones.mcones.mcone.isMCone.mcone_M_sep}).
    \item[$(\mathrm{Msnorm})$, dual-norm-sup] for every $x \ne 0$ in
      $C$ and every $\varepsilon > 0$, there exists
      $m \in \Mtest_{\arzero}$ such that
      $\cnorm{x} \le m\ \arzero\ x + \varepsilon$
      (\rocq{Icones.mcones.mcone.isMCone.mcone_M_norm}).
  \end{description}
  The Rocq mixin is registered against a \texttt{coneType\ R} via
  the HB structure \texttt{MCone.type\ Ar} (alias
  \texttt{mconeType\ Ar}).
\end{definition}

\begin{remark}
  Two encoding choices deserve emphasis.

  \emph{First, the form of $(\mathrm{Msnorm})$.} The paper writes the
  axiom in the form $\cnorm{x} \le m(x) / \cnorm{m} + \varepsilon$.
  Every test in our encoding already satisfies the operator-norm
  bound $\cnorm{m} \le 1$ via the field
  \rocq{Icones.mcones.mcone.test_norm_le}, so the simplified form
  $\cnorm{x} \le m(x) + \varepsilon$ adopted in
  \texttt{isMCone.mcone\_M\_norm} is logically equivalent: the
  $/\cnorm{m}$ factor would only weaken the bound, and we do not
  formalise the dual norm as a separate notion at the level of
  individual tests.

  \emph{Second, tests are real-valued and not bundled as cone
  morphisms.} A test $m\ r$ is morally an element of
  $\Cones(C, \mathbf{1})$ for each fixed $r$, but we keep the
  pointwise linearity and $\omega$-continuity facts inline rather
  than re-bundling each test as a \texttt{cones\_hom}; this avoids
  weaving the $\Rnn$/$\{nonneg\ R\}$-cast machinery into every test
  manipulation.
\end{remark}

\begin{definition}[Paper Def 3.6, measurable cone]\label{def:mcone}
  \uses{def:isMCone}
  \rocq{Icones.mcones.mcone.MCone}\rocqok
  A \emph{measurable cone} is a cone $C$ equipped with a measurability
  structure $\Mtest$. The HB structure
  $\MCone.type\ \Ar$ packages a cone together with an
  \texttt{isMCone} mixin; we write $\MCone$ for an inhabitant of this
  structure.
\end{definition}

%% Measurable paths. --------------------------------------------------

\begin{definition}[Paper Def 3.7, measurable path]\label{def:meas-path}
  \uses{def:mcone}
  \rocq{Icones.mcones.mcone.is_measurable_path}\rocqok
  Let $C$ be a measurable cone and $X \in \Ar$. A \emph{measurable
  path of arity $X$} in $C$ is a function $\gamma : \ar\_carrier X
  \to C$ that is
  \begin{itemize}
    \item norm-bounded: there exists $M : \R$ with $\cnorm{\gamma\ r}
      \le M$ for every $r$;
    \item test-measurable: for every $Y \in \Ar$ and every
      $m \in \Mtest_Y$, the function
      $(s, r) \mapsto m\ s\ (\gamma\ r)$ on
      $\ar\_carrier Y \times \ar\_carrier X$ is measurable to $\R$.
  \end{itemize}
\end{definition}

\begin{lemma}[Paper Lemma 3.9, constant paths]\label{lem:const-path}
  \uses{def:meas-path}
  \rocq{Icones.mcones.mcone.const_path_measurable}\rocqok
  For every $x \in C$, the constant function
  $\gamma : r \mapsto x$ is a measurable path of arity $X$.
\end{lemma}
\begin{proof}
  The boundedness witness is $M = \cnorm{x}$. For
  test-measurability, the test field
  \rocq{Icones.mcones.mcone.test_meas} gives measurability of
  $r \mapsto m\ r\ x$ when $\cnorm{x} \le 1$; we precompose with the
  first projection. If $\cnorm{x} > 1$ we rescale $x$ to
  $x' = (\cnorm{x})^{-1} \cdot x$ (which lies in the unit ball) and
  use the test's $\Rnn$-linearity in $x$ to recover the result by
  multiplying by the constant $\cnorm{x}$.
\end{proof}

\begin{lemma}[Paper Lemma 3.10, reindexing of paths]%
  \label{lem:path-reindex}
  \uses{def:meas-path}
  \rocq{Icones.mcones.mcone.reindex_path_measurable}\rocqok
  Let $\gamma : X \to C$ be a measurable path of arity $X$ and let
  $\varphi \in \Ar(Y, X)$. Then $\gamma \circ \varphi$ is a
  measurable path of arity $Y$.
\end{lemma}
\begin{proof}
  Boundedness is immediate (same $M$). For measurability, given
  $m \in \Mtest_{Y'}$ the map
  $(s', s) \mapsto m\ s'\ (\gamma\ (\varphi\ s))$ factors as
  $(s', r) \mapsto m\ s'\ (\gamma\ r)$ (measurable by hypothesis on
  $\gamma$) precomposed with the pairing $(s', s) \mapsto (s',
  \varphi\ s)$, itself measurable as a pair of measurable
  components.
\end{proof}

%%%%%%%%%%%%%%%%%%%%%%%%%%%%%%%%%%%%%%%%%%%%%%%%%%%%%%%%%%%%%%%%%%%%%%%%
\section{The category \texorpdfstring{$\MCones$}{MCones}}%
\label{sec:cat-mcones}

\begin{definition}[Paper Def 3.13, the category $\MCones$]%
  \label{def:mcones-hom}
  \uses{def:mcone, def:meas-path}
  \rocq{Icones.mcones.mcone_cat.mcones_hom}\rocqok
  Objects of $\MCones$ are measurable cones. A morphism in
  $\MCones(B, C)$ is a $\Cones$-morphism $f : B \to C$ (linear,
  $\omega$-continuous, norm-non-expansive) together with a
  path-preservation property: for every $X \in \Ar$ and every
  measurable path $\gamma : X \to B$, the composition
  $f \circ \gamma$ is a measurable path $X \to C$
  (\rocq{Icones.mcones.mcone_cat.mcones_hom_pres_path}).
  The identity and composition are inherited from $\Cones$ and the
  path-preservation field is preserved
  (\rocq{Icones.mcones.mcone_cat.mcones_id},
   \rocq{Icones.mcones.mcone_cat.mcones_comp},
   \rocq{Icones.mcones.mcone_cat.mcones_compIl},
   \rocq{Icones.mcones.mcone_cat.mcones_compIr},
   \rocq{Icones.mcones.mcone_cat.mcones_compA}).\rocqok
\end{definition}

\begin{proposition}[Paper Prop 3.11, dual-norm sup characterisation]%
  \label{prop:mcone-norm-sup}
  \uses{def:mcone}
  \rocq{Icones.mcones.mcone_cat.mcone_norm_le_pairing_ub}\rocqok
  Let $B$ be a measurable cone. For every $x \in B$,
  \[
    \cnorm{x} = \sup \bigl\{ m\ \arzero\ x \bigm| m \in \Mtest_{\arzero} \bigr\}.
  \]
\end{proposition}
\begin{proof}
  The $\ge$ direction is the operator-norm bound
  $m\ r\ x \le \cnorm{x}$ from
  Definition~\ref{def:test-of}
  (\rocq{Icones.mcones.mcone_cat.mcone_test_pairing_ub}).
  For $\le$, suppose $x \ne 0$ and $M$ is an upper bound of the
  pairing set. By $(\mathrm{Msnorm})$, for every $\varepsilon > 0$
  there exists $m \in \Mtest_{\arzero}$ with
  $\cnorm{x} \le m\ \arzero\ x + \varepsilon \le M + \varepsilon$
  (\rocq{Icones.mcones.mcone_cat.mcone_test_pairing_adherent});
  letting $\varepsilon \to 0$ yields $\cnorm{x} \le M$. The case
  $x = 0$ is immediate from $\cnorm{0} = 0$.
\end{proof}

\begin{remark}[Paper Rem 3.12]
  The point of $(\mathrm{Msnorm})$ is precisely
  Proposition~\ref{prop:mcone-norm-sup}: in the absence of a
  Hahn--Banach theorem for cones (the paper uses Remark~2.16 as a
  counter-example), separation in the operator-norm sense is taken
  as an \emph{axiom} on the tests rather than derived. This axiom
  survives all the limit constructions that follow in \S4--\S6.
\end{remark}

%% Rescaling. ---------------------------------------------------------

\begin{definition}[Paper Def 3.15, the rescaled cone $\alpha B$]%
  \label{def:alpha-rescale}
  \uses{def:mcone}
  \rocq{Icones.mcones.mcone_cat.alpha_rescale_car}\rocqok
  Let $B$ be a measurable cone and $\alpha > 0$. The \emph{rescaling}
  $\alpha B$ is the measurable cone with the same underlying carrier
  (wrapped in a thin record \texttt{alpha\_rescale\_car} to give HB a
  distinct identity), the same precone operations, the rescaled norm
  \[
    \cnorm{x}_{\alpha B} = \alpha^{-1} \cnorm{x}_B,
  \]
  and the test family obtained by scaling every test of $B$ by the
  $\Rnn$-coefficient $\alpha^{-1}$ (so that the operator-norm bound
  $m\ r\ x \le \cnorm{x}_{\alpha B}$ holds for the rescaled norm).
  The full \texttt{isPrecone}/\texttt{isCone}/\texttt{isMCone} HB
  tower is registered on
  \texttt{alpha\_rescale\_car \(\alpha\) B} in
  \texttt{theories/mcones/mcone\_cat.v}.\rocqok
  We have $B_{\alpha B} = \alpha B_B = \{ x \in B \mid \cnorm{x}_B
  \le \alpha \}$.
\end{definition}

%%%%%%%%%%%%%%%%%%%%%%%%%%%%%%%%%%%%%%%%%%%%%%%%%%%%%%%%%%%%%%%%%%%%%%%%
\section{Example: finite measures
  \texorpdfstring{$\FMeas(X)$}{FMeas(X)}}\label{sec:fmeas}

Let $X$ be a measurable space (not necessarily in $\Ar$). The set
$\FMeas(X)$ of finite measures on $X$ has been equipped with a cone
structure in \S\ref{ch:cones}; this section endows it with a
measurability structure in the sense of
Definition~\ref{def:isMCone}.

\begin{definition}[Paper §2.2, the carrier $\FMeas(X)$]%
  \label{def:fmeas-carrier}
  \uses{def:meas-subcat}
  \rocq{Icones.mcones.fmeas.fmeas}\rocqok
  In Rocq, an element of \texttt{fmeas R X} is a record wrapping a
  measure $\mu : \{\text{measure}\ \mathsf{set}\ X \to \Rext\}$
  together with two structural invariants:
  \begin{itemize}
    \item finiteness: $\mu(U) \in \R$ for every measurable $U$
      (\rocq{Icones.mcones.fmeas.fmeas_finP});
    \item canonical extension: $\mu(U) = 0$ for every non-measurable
      $U$ (\rocq{Icones.mcones.fmeas.fmeas_canon}).
  \end{itemize}
  The canonicality invariant lets equality of \texttt{fmeas R X}
  elements reduce to equality on the $\sigma$-algebra
  (\rocq{Icones.mcones.fmeas.fmeas_eq}).
\end{definition}

The full HB tower
($\PCone \to \Cone \to \MCone$) on \texttt{fmeas R X} is registered in
\texttt{theories/mcones/fmeas.v}. We focus here on the measurability
structure, which uses the simplest test family of the paper
(\S3.2.1).

\begin{definition}[Paper §3.2.1, the tests $e_U$ on $\FMeas(X)$]%
  \label{def:fmeas-tests}
  \uses{def:fmeas-carrier, def:isMCone}
  \rocq{Icones.mcones.fmeas.isMCone}\rocqok
  For each measurable $U \in \sigma_X$, define the test
  $\evtest{U} : \ar\_carrier Y \times \FMeas(X) \to \Rnn$ by
  $\evtest{U}(s, \mu) = \mu(U)$ (constant in $s$). The family
  $\Mtest_Y = \{ \evtest{U} \mid U \in \sigma_X \}$ is a measurability
  structure on $\FMeas(X)$: measurability is trivial (the partial
  application $r \mapsto \mu(U)$ is constant), linearity and
  $\omega$-continuity in $\mu$ come from the underlying measure
  algebra, the reindexing closure $(\mathrm{Mscomp})$ is exact (the
  reindexed test is again of the form $\evtest{U}$), separation
  $(\mathrm{Mssep})$ is the fact that two finite measures agreeing on
  every measurable set are equal, and $(\mathrm{Msnorm})$ follows
  from $\cnorm{\mu} = \mu(X)$ together with adherence of the
  measure to its total mass.
\end{definition}

\begin{lemma}[Paper Lemma 3.17, pushforward functor]%
  \label{lem:fmeas-push}
  \uses{def:fmeas-tests, def:mcones-hom}
  \rocq{Icones.mcones.fmeas.fmeas_push}\rocqok
  For every $\varphi \in \Ar(X, Y)$ the pushforward
  $\mu \mapsto \varphi_*\mu$ defined by
  $(\varphi_*\mu)(V) = \mu(\varphi^{-1}(V))$ is a morphism
  $\FMeas(X) \to \FMeas(Y)$ in $\MCones$. The assignment
  $X \mapsto \FMeas(X)$, $\varphi \mapsto \varphi_*$ is functorial
  (\rocq{Icones.mcones.fmeas.fmeas_push_id},
  \rocq{Icones.mcones.fmeas.fmeas_push_comp}), giving a functor
  $\FMeas : \Ar \to \MCones$.
\end{lemma}
\begin{proof}
  Linearity, $\omega$-continuity and the norm bound
  $\cnorm{\varphi_*\mu} \le \cnorm{\mu}$ come from the pointwise
  definition of the pushforward and the cone structure of $\FMeas$.
  For measurability of paths: given a measurable path
  $\kappa : Y' \to \FMeas(X)$ and a test $\evtest{V}$ on $\FMeas(Y)$
  with $V \in \sigma_Y$, the function
  $(s'', s') \mapsto \evtest{V}(s'', (\varphi_* \kappa(s')))
   = \kappa(s')(\varphi^{-1}(V))$
  is measurable because $\kappa$ is a measurable path and
  $\varphi^{-1}(V) \in \sigma_X$.
\end{proof}

\begin{remark}
  The Rocq encoding uses a finiteness + canonicality invariant
  wrapper \texttt{fmeas R X} around the mathcomp-analysis bundle
  \texttt{\{measure ... -> \(\overline{\R}\)\}}. The wrapper is what
  makes $\sigma$-algebra-pointwise equality
  (\rocq{Icones.mcones.fmeas.fmeas_eq}) a usable equality predicate,
  which in turn makes every cone-axiom proof on \texttt{fmeas} clean.
\end{remark}

%%%%%%%%%%%%%%%%%%%%%%%%%%%%%%%%%%%%%%%%%%%%%%%%%%%%%%%%%%%%%%%%%%%%%%%%
\section{Example: paths
  \texorpdfstring{$\Path(X, B)$}{Path(X, B)}}\label{sec:path}

Let $B$ be a measurable cone and $X \in \Ar$. The set of measurable
paths $X \to B$ (Definition~\ref{def:meas-path}) is itself a measurable
cone, with pointwise algebra, sup-norm, and the test family $\varphi
\trianglerhd m$ described below.

\begin{definition}[Paper §3.2.2, the carrier of $\Path(X, B)$]%
  \label{def:path-carrier}
  \uses{def:meas-path}
  \rocq{Icones.mcones.path.path_car}\rocqok
  An element of \texttt{path\_car Ar X B} is a record wrapping a
  function $\gamma : \ar\_carrier X \to B$ together with the
  measurable-path proof
  \texttt{is\_measurable\_path \(\gamma\)}.
  Two paths are equal iff their underlying functions agree pointwise
  (\rocq{Icones.mcones.path.path_eq}).
\end{definition}

The pointwise algebraic structure on \texttt{path\_car Ar X B} is
inherited from $B$: addition, scalar multiplication and the order are
defined pointwise, the boundedness witness of the result combines the
witnesses of the operands, and test-measurability of the result is
inherited from pointwise linearity and $\omega$-continuity of each
test in its $B$-argument. The full HB tower is registered in
\texttt{theories/mcones/path.v}\rocqok:
\begin{itemize}
  \item the precone instance
    (\rocq{Icones.mcones.path.path_addA},
    \rocq{Icones.mcones.path.path_addC},
    \rocq{Icones.mcones.path.path_scale_A});
  \item the norm $\cnorm{\gamma} = \sup_r \cnorm{\gamma(r)}$
    (\rocq{Icones.mcones.path.path_norm});
  \item the cone instance with the $(\mathrm{Normc})$ witness
    $\bigsup_n \gamma_n$ built pointwise from $B$'s $\cone\_sup\_ball$
    (\rocq{Icones.mcones.path.path_sup_ball},
    \rocq{Icones.mcones.path.path_sup_ball_meas}). The
    $(\mathrm{Normc})$-measurability proof uses monotone convergence
    on the increasing chain of measurable functions
    $(s, r) \mapsto m\ s\ (\gamma_n\ r)$.
\end{itemize}

\begin{definition}[Paper §3.2.2, the test family $\varphi
  \trianglerhd m$]\label{def:path-test}
  \uses{def:test-of, def:path-carrier}
  \rocq{Icones.mcones.path.path_test}\rocqok
  For $Y \in \Ar$, $\varphi \in \Ar(Y, X)$ and a test $m$ of arity
  $Y$ on $B$, define the test $\varphi \trianglerhd m$ of arity $Y$
  on $\Path(X, B)$ by
  \[
    (\varphi \trianglerhd m)(s, \gamma) = m(s, \gamma(\varphi(s))) .
  \]
  Linearity, $\omega$-continuity and the norm bound for $\varphi
  \trianglerhd m$ all reduce to the corresponding properties of $m$.
  The family
  \[
    \Mtest_Y = \{ \varphi \trianglerhd m \mid
      \varphi \in \Ar(Y, X),\ m \in \Mtest_Y \text{ on } B \}
  \]
  is a measurability structure on $\Path(X, B)$
  (\rocq{Icones.mcones.path.isMCone}).\rocqok
\end{definition}

\begin{remark}
  The four axioms of Definition~\ref{def:isMCone} for the family
  $\varphi \trianglerhd m$:
  \begin{description}
    \item[$(\mathrm{Mscomp})$] For
      $\psi \in \Ar(Y', Y)$,
      $(\varphi \trianglerhd m) \circ (\psi \times P)
        = (\varphi \circ \psi) \trianglerhd (\reindex{\psi}{m})$
      (\rocq{Icones.mcones.path.path_mcone_M_comp}).
    \item[$(\mathrm{Mssep})$] If $\gamma_1, \gamma_2 \in \Path(X, B)$
      satisfy $p(\gamma_1) = p(\gamma_2)$ for every
      $p \in \Mtest_{\arzero}$, then for each $r \in \ar\_carrier X$
      the constant path $\arzero \to X$ with value $r$ (encoded by
      \texttt{const\_r r}) gives
      $m(\gamma_1(r)) = m(\gamma_2(r))$ for every
      $m \in \Mtest_{\arzero}^B$, and $(\mathrm{Mssep})$ in $B$
      finishes the argument
      (\rocq{Icones.mcones.path.path_mcone_M_sep}).
    \item[$(\mathrm{Msnorm})$] For $\gamma \ne 0$ and
      $\varepsilon > 0$, sup-adherence on
      $\{\cnorm{\gamma(r)} \mid r\}$ yields $r_0$ with
      $\cnorm{\gamma} \le \cnorm{\gamma(r_0)} + \varepsilon/2$;
      applying $(\mathrm{Msnorm})$ of $B$ at $\gamma(r_0)$ (if
      non-zero) gives a test $m$ with the matching bound, and the
      witness on $\Path(X, B)$ is
      $(\textsf{const}_{r_0}) \trianglerhd m$
      (\rocq{Icones.mcones.path.path_mcone_M_norm}). The degenerate
      case $\gamma(r_0) = 0$ is handled by picking some other $r_1$
      where $\gamma(r_1) \ne 0$.
  \end{description}
\end{remark}

%% Flattening iso. ----------------------------------------------------

\begin{lemma}[Paper Lemma 3.19, flattening iso]\label{lem:path-flatten}
  \uses{def:path-carrier, def:mcones-hom, lem:ar-prod-up}
  \rocq{Icones.mcones.path.path_fl_fun}
  \rocq{Icones.mcones.path.path_fl_inv_fun}
  \rocq{Icones.mcones.path.path_fl_fun_inv}\rocqok
  For $X, Y \in \Ar$ and $B$ a measurable cone there is, at the level
  of underlying bivariate functions, an iso
  \[
    \mathrm{fl}_{X, Y} : \Path(X, \Path(Y, B))
    \xrightarrow{\sim}
    \Path(X \times Y, B),
    \qquad
    \mathrm{fl}_{X, Y}(\eta)(r, s) = \eta(r)(s).
  \]
  The inverse map sends a path $\eta : X \times Y \to B$ to
  $r \mapsto (s \mapsto \eta(r, s))$, and the two are mutual
  inverses. Consequently the parameter-swap
  $\mathrm{fl}_{Y, X}^{-1} \circ \mathrm{fl}_{X, Y}$ is an iso
  $\Path(X, \Path(Y, B)) \xrightarrow{\sim} \Path(Y, \Path(X, B))$.
\end{lemma}
\begin{proof}[Proof sketch]
  Boundedness is immediate ($\cnorm{\mathrm{fl}(\eta)} = \cnorm{\eta}$
  as sets of pointwise norms). For measurability of
  $\mathrm{fl}(\eta) : X \times Y \to B$: given a test $m$ on $B$
  with arity $Y'$, the function
  $(s', r, s) \mapsto m(s', \eta(r)(s))$ rewrites as
  $\pi_2 \trianglerhd m'$ applied to $(s, \eta(r))$ with
  $m' = m \circ (\pi_1 \times B)$, which is measurable by the
  measurability of $\eta$ as a path in $\Path(Y, B)$. The reverse
  direction is symmetric, using the test $\varphi \trianglerhd m$ on
  $\Path(X \times Y, B)$ together with the universal property of
  $\ar\_prod$ (Lemma~\ref{lem:ar-prod-up}).
\end{proof}

\begin{remark}[Encoding of Lemma 3.19]
  In the paper, Lemma~3.19 is phrased as an iso in $\MCones$. In
  Rocq, the carrier of $\ar\_prod\ X\ Y$ is only \emph{propositionally}
  equal to $\ar\_carrier X \times \ar\_carrier Y$ (the
  \texttt{ar\_prod\_carrier\_eq} field of \texttt{MeasSubcat}), so
  rephrasing $\mathrm{fl}$ as a function
  $\Path(X, \Path(Y, B)) \to \Path(\ar\_prod X Y, B)$ requires
  transporting along this equation. The function-level content
  (definitional pointwise inverse) is fully formalised, as is the
  norm equality at the bivariate-function level. The full
  \texttt{cones\_hom}-packaging of $\mathrm{fl}$ uses the M3 wave 7
  cast-measurability fields of $\ar\_prod$
  (\rocq{Icones.mcones.ar.ar_prod_cast_meas},
  \rocq{Icones.mcones.ar.ar_prod_uncast_meas}); the lifting is a
  routine but voluminous transport, deferred to the tensor-product
  layer.\notready
\end{remark}

%%%%%%%%%%%%%%%%%%%%%%%%%%%%%%%%%%%%%%%%%%%%%%%%%%%%%%%%%%%%%%%%%%%%%%%%
\section{Design choices}\label{sec:mcones-design}

\paragraph{Why parametrise by $\Ar$.}
Paper Remark~3.1 notes that taking $\Ar = \{\R^n \mid n \ge 0\}$ would
already be expressive enough for everything that follows. We keep
$\Ar$ a parameter for three reasons. First, it is paper-faithful:
\S3 introduces $\Ar$ as a parameter, the integrability theory of
\S4--\S6 never uses anything beyond the closure properties stated in
Definition~\ref{def:meas-subcat}, and reproducing this dependency
clearly in Rocq makes the link between the formalisation and the
paper transparent. Second, downstream cones (notably $\Path(X, B)$ in
this chapter) are themselves parametrised by an object of $\Ar$, and
threading $\Ar$ as a uniform parameter keeps the HB structure-builder
happy. Third, future extensions to other categories of measurable
spaces (Polish, standard Borel, $\mathrm{QBS}$, $\dots$) can be done
just by instantiating $\Ar$.

\paragraph{Propositional vs definitional cartesian product.}
The cleanest mathematical statement is that $\ar\_prod$ is
\emph{definitionally} the mathcomp-analysis product
$\mathsf{measurableType}$; we chose a propositional equality
\texttt{ar\_prod\_carrier\_eq} instead because we wanted the $\Ar$
parameter to be an abstract \texttt{Record} rather than a concrete
constructor. The price is that downstream code occasionally has to
transport across the equality; the cast/uncast measurability fields
(M3 wave 7) absorb this overhead at the abstraction boundary.

\paragraph{The test-family representation.}
A test is a \texttt{Record\ test\_of} carrying (a) the underlying
function, (b) the measurability proof, (c) the bounds
$0 \le m\ r\ x \le 1$ on the unit ball, (d) the
linearity-in-the-second-argument proofs, (e) the $\omega$-continuity
proof, and (f) the operator-norm bound. Bundling the proofs with the
function (rather than via separate predicates) avoids re-exposing the
same hypotheses at every use site, and gives a clean
\emph{equational} reasoning principle:
\rocq{Icones.mcones.mcone.test_eq} says two tests are equal as soon
as their underlying functions agree pointwise.

\paragraph{Why $\fmeas$ wraps measures.}
A \texttt{fmeas R X} is a record around a mathcomp-analysis
\texttt{\{measure ...\}} with two invariants: finiteness on every
measurable set, and canonical extension by $0$ off the
$\sigma$-algebra. The canonicality invariant is the technical heart:
it makes equality on the $\sigma$-algebra the right equality
predicate for elements of $\fmeas$, which in turn is what every cone
axiom on $\fmeas$ needs. Without it, two finite measures could agree
on every measurable set but differ on non-measurable sets, and one
would be forced to use setoids throughout the cone-axiom proofs.
